% Chapter 5 -- Orthogonality, projections, and QR
\chapter{Orthogonality, projections, and QR}\label{ch:orthogonality}

Chapter~3 built the four fundamental subspaces; this chapter gives them
\emph{geometry}. Once you have an inner product you can measure angles and
lengths, split a vector into ``part inside a subspace'' and ``part outside,''
and replace any basis by an orthonormal one that makes every coordinate
computation a dot product. These are the tools behind least squares (the
next chapter) and behind almost every numerically stable algorithm in this
book. After this chapter you will be able to project onto lines and
subspaces, run Gram--Schmidt by hand and in code, factor $A=QR$, and see
why classical Gram--Schmidt loses orthogonality while the modified version
does not.

\section{Inner product, norm, angle}\label{sec:orth-inner}

Coordinates tell you where a vector sits relative to axes. They do not, by
themselves, tell you how long it is or what angle two vectors make --- for
that you need one more piece of structure.

\begin{definition}[Inner product]\index{inner product}\index{dot product}
The \emph{(standard) inner product}, or \emph{dot product}, of
$\mathbf u=(u_1,\dots,u_n)$ and $\mathbf v=(v_1,\dots,v_n)\in\mathbb R^n$ is
\[
\langle\mathbf u,\mathbf v\rangle=\mathbf u^{T}\mathbf v
=u_1v_1+u_2v_2+\cdots+u_nv_n .
\]
It is symmetric ($\langle\mathbf u,\mathbf v\rangle=\langle\mathbf v,\mathbf u\rangle$),
linear in each argument, and positive definite:
$\langle\mathbf v,\mathbf v\rangle=\|\mathbf v\|^2>0$ unless $\mathbf v=\mathbf 0$.
\end{definition}

\begin{definition}[Norm and angle]\index{norm}\index{angle between vectors}
$\|\mathbf v\|=\sqrt{\langle\mathbf v,\mathbf v\rangle}$ is the \emph{norm} (length)
of $\mathbf v$, and the angle between two nonzero vectors satisfies
\[
\cos\theta=\frac{\langle\mathbf u,\mathbf v\rangle}{\|\mathbf u\|\,\|\mathbf v\|},
\qquad 0\le\theta\le\pi .
\]
\end{definition}

Two facts make the angle formula usable rather than decorative:

\begin{theorem}[Cauchy--Schwarz]\index{Cauchy--Schwarz inequality}
$|\langle\mathbf u,\mathbf v\rangle|\le\|\mathbf u\|\,\|\mathbf v\|$, with equality
exactly when $\mathbf u,\mathbf v$ are parallel.
\end{theorem}

\begin{theorem}[Pythagoras]\index{Pythagoras}
If $\langle\mathbf u,\mathbf v\rangle=0$ then
$\|\mathbf u+\mathbf v\|^2=\|\mathbf u\|^2+\|\mathbf v\|^2$.
\end{theorem}

\begin{proof}[Proof of Pythagoras]
$\|\mathbf u+\mathbf v\|^2=\langle\mathbf u+\mathbf v,\mathbf u+\mathbf v\rangle
=\|\mathbf u\|^2+2\langle\mathbf u,\mathbf v\rangle+\|\mathbf v\|^2$, and the
cross term vanishes by hypothesis.
\end{proof}

\begin{definition}[Orthogonal vectors]\index{orthogonal vectors}
$\mathbf u,\mathbf v$ are \emph{orthogonal}, written $\mathbf u\perp\mathbf v$,
when $\langle\mathbf u,\mathbf v\rangle=0$. A vector is a \emph{unit vector}
when $\|\mathbf v\|=1$; any nonzero $\mathbf v$ has unit vector
$\hat{\mathbf v}=\mathbf v/\|\mathbf v\|$ in the same direction.
\end{definition}

\begin{exercise}[easy]\label{ex:orth-1}
For $\mathbf u=(1,2,-1)$, $\mathbf v=(3,-1,1)$, compute $\langle\mathbf u,\mathbf v\rangle$,
$\|\mathbf u\|$, and $\cos\theta$.
\begin{solution}
$\langle\mathbf u,\mathbf v\rangle=3-2-1=0$, so the vectors are orthogonal and
$\theta=\pi/2$. $\|\mathbf u\|=\sqrt{1+4+1}=\sqrt6$.
\end{solution}
\end{exercise}

\begin{example}[Angle and cosine similarity]\label{ex:cossim}
Find the angle between $\mathbf{a}=(1,2,2)$ and
$\mathbf{b}=(2,-1,2)$, and their \emph{cosine similarity}
$\cos\theta$ --- the statistic recommendation engines actually
use.
\par\smallskip
\textbf{Strategy.} $\cos\theta=\mathbf a\cdot\mathbf b/
(\|\mathbf a\|\|\mathbf b\|)$; all three numbers are one dot product
each. Compare also $\|\mathbf a-\mathbf b\|$: distance and angle answer
different questions.
\begin{solution}
$\mathbf a\cdot\mathbf b=2-2+4=4$; $\|\mathbf a\|=3$,
$\|\mathbf b\|=3$; so $\cos\theta=4/9\approx0.444$ and
$\theta\approx63.6^\circ$. Meanwhile $\|\mathbf a-\mathbf
b\|=\|(-1,3,0)\|=\sqrt{10}\approx3.16$: vectors nearly a right angle
apart can still be close in raw distance when they are short. For
sparse high-dimensional taste vectors the \emph{angle} is the honest
similarity --- magnitude is mostly noise there. \checkmark
\end{solution}
\end{example}

\section{Orthogonal and orthonormal sets}\label{sec:orth-sets}

\begin{definition}[Orthogonal / orthonormal set]\index{orthonormal set}
Vectors $\mathbf q_1,\dots,\mathbf q_k$ form an \emph{orthogonal set} when
$\mathbf q_i\perp\mathbf q_j$ for $i\ne j$. If additionally each is a unit
vector, the set is \emph{orthonormal}: $\mathbf q_i^{T}\mathbf q_j=\delta_{ij}$.
An orthonormal set that is also a basis is an \emph{orthonormal basis}.
\end{definition}

\begin{theorem}\label{th:orth-indep}
A finite orthogonal set of nonzero vectors is linearly independent.
\end{theorem}

\begin{proof}
Suppose $c_1\mathbf q_1+\cdots+c_k\mathbf q_k=\mathbf0$. Take the inner product
with $\mathbf q_j$: every term except $c_j\|\mathbf q_j\|^2$ is orthogonal to
$\mathbf q_j$, so $c_j\|\mathbf q_j\|^2=0$ and hence $c_j=0$ (each
$\|\mathbf q_j\|\ne0$). This holds for every $j$, so all coefficients vanish.
\end{proof}

The payoff: in an orthonormal basis, coordinates are dot products.

\begin{theorem}\label{th:orth-coords}
If $\mathcal Q=\{\mathbf q_1,\dots,\mathbf q_n\}$ is an orthonormal basis of
$\mathbb R^n$ and $\mathbf v=c_1\mathbf q_1+\cdots+c_n\mathbf q_n$, then
$c_j=\mathbf q_j^{T}\mathbf v$.
\end{theorem}

Compare with a general basis (ch.~3): extracting $[\mathbf v]_{\mathcal B}$
meant solving a linear system. With an orthonormal basis it is $n$ dot
products --- and each coefficient is computed \emph{independently}, so an
error in one cannot corrupt the others.

\begin{definition}[Orthogonal matrix]\index{orthogonal matrix}
A square matrix $Q$ is \emph{orthogonal} when $Q^{T}Q=I$, i.e.
$Q^{-1}=Q^{T}$: its columns form an orthonormal basis of $\mathbb R^n$.
\end{definition}

Orthogonal matrices preserve everything geometric:
$\|Q\mathbf v\|=\|\mathbf v\|$ and
\[
\langle Q\mathbf u,Q\mathbf v\rangle
=\mathbf u^{T}Q^{T}Q\mathbf v=\langle\mathbf u,\mathbf v\rangle .
\]
Rotations and reflections from ch.~4 are exactly the 2D orthogonal
matrices ($\det Q=\pm1$). Their numerical value is just as large:
multiplying by $Q$ or $Q^{T}$ cannot amplify rounding error, which is
why stable algorithms are built out of orthogonal factors.

\begin{figure}[htbp]\centering
% Figure 5.x: an orthonormal frame is a rotated standard frame.
\begin{tikzpicture}[scale=1.15,>=stealth]
  % standard frame (faint)
  \draw[->,very thick,gray!50] (0,0) -- (2.2,0) node[below,font=\small] {$\mathbf e_1$};
  \draw[->,very thick,gray!50] (0,0) -- (0,2.2) node[left,font=\small] {$\mathbf e_2$};
  % orthonormal frame rotated ~30 deg
  \draw[->,very thick,accent] (0,0) -- (1.90,1.10) node[right,font=\small] {$\mathbf q_1=(\cos30^\circ,\sin30^\circ)$};
  \draw[->,very thick,accent] (0,0) -- (-1.10,1.90) node[above,font=\small] {$\mathbf q_2=(-\sin30^\circ,\cos30^\circ)$};
  % right angle marker between q1 q2
  \draw (0.45,0.26) -- (0.19,0.71) -- (-0.26,0.45);
  % unit circle hint
  \draw[dotted,gray] (0,0) circle (2.2);
  \node[font=\footnotesize,align=center] at (3.6,-0.6) {$\mathbf q_1\cdot\mathbf q_2=0$,\ \ $\|\mathbf q_i\|=1$:\\ orthonormal = rotated,\\ never stretched};
\end{tikzpicture}

\caption{An orthonormal pair is the standard frame, rotated. Notice
what orthonormality buys: the coordinates of a vector in the
$\mathbf q$-frame are just the two projections --- no solving needed.}
\label{fig:frame}
\end{figure}

\section{Projection onto a line and onto a subspace}\label{sec:orth-proj}

The concrete problem: given a vector $\mathbf b$ and a subspace $W$, which
point of $W$ is closest to $\mathbf b$? It is the foot of the perpendicular
--- the \emph{projection}.

The derivation itself is worth doing once properly
(Section~\ref{sec:handderiv5}); the result is
\[
\mathbf p=\frac{\mathbf a^{T}\mathbf b}{\mathbf a^{T}\mathbf a}\,\mathbf a,
\qquad
P=\frac{\mathbf a\mathbf a^{T}}{\mathbf a^{T}\mathbf a}
\]
--- a rank-1 matrix satisfying $P^2=P$ and $P^{T}=P$.

\begin{figure}[htbp]
\centering
\begin{tikzpicture}[scale=1.1, >=stealth]
  \draw[->, gray] (-0.5,0)--(4.5,0);
  \draw[->, gray] (0,-0.5)--(0,3);
  \draw[accent, very thick] (0,0)--(4,1) node[right]{$\mathbf a$};
  \draw[->, ink, very thick] (0,0)--(2.4,2.6) node[above]{$\mathbf b$};
  \draw[->, good, very thick] (0,0)--(2.08,0.52) node[below]{$\mathbf p=\hat b\,\mathbf a$};
  \draw[->, amber, thick] (2.08,0.52)--(2.4,2.6) node[midway,right]{$\mathbf e=\mathbf b-\mathbf p$};
  \draw (2.08,0.52) +(0,0.18) -- +(0.13,0.18) -- +(0.13,0);
  \node[align=left] at (3.2,2.5) {\small residual $\perp$ line:\\ \small $\mathbf e\cdot\mathbf a=0$};
\end{tikzpicture}

\caption{Projection of $\mathbf b$ onto the line through $\mathbf a$. Notice
that the formula follows entirely from one requirement: the residual
$\mathbf e=\mathbf b-\mathbf p$ is perpendicular to the line. Regenerate:
\ttfamily figures/ch05\_projection\_line.tex.}
\label{fig:proj-line}
\end{figure}

The same ``residual is perpendicular'' condition works for any subspace.
If $\mathbf q_1,\dots,\mathbf q_k$ is an \emph{orthonormal} basis of $W$, the
projection splits into independent one-dimensional pieces:

\begin{proposition}\label{th:proj-onb}
For orthonormal $\mathbf q_1,\dots,\mathbf q_k$ spanning $W$,
\[
\mathbf p=(\mathbf q_1^{T}\mathbf b)\,\mathbf q_1+\cdots+
(\mathbf q_k^{T}\mathbf b)\,\mathbf q_k .
\]
\end{proposition}

\begin{proof}
Let $\mathbf p$ have that form. Then $\mathbf p\in W$, and for each basis
vector, $\mathbf q_j^{T}(\mathbf b-\mathbf p)
=\mathbf q_j^{T}\mathbf b-\mathbf q_j^{T}\mathbf b=0$, so the residual is
orthogonal to every basis vector, hence to all of $W$.
\end{proof}

For a \emph{general} basis $A=[\mathbf a_1\ \cdots\ \mathbf a_k]$ of $W$
(tall, independent columns), the condition $A^{T}(\mathbf b-A\mathbf x)=\mathbf0$
gives the normal equations $A^{T}A\mathbf x=A^{T}\mathbf b$, and

\begin{theorem}\label{th:proj-matrix}
If $A$ has linearly independent columns, the orthogonal projection onto
$\operatorname{col}(A)$ is
\[
P=A(A^{T}A)^{-1}A^{T},\qquad P^2=P,\quad P^{T}=P .
\]
\end{theorem}

\begin{tripwire}[{Believing $A^{T}A$ is always invertible}]
$A^{T}A$ is invertible \emph{if and only if} $A$'s columns are linearly
independent (ch.~8 proves it). Feed it a rank-deficient $A$ and the
inverse does not exist --- which is exactly what least squares does when
your data has redundant features.
\end{tripwire}

\begin{example}[Projection onto a line and a plane]\label{ex:proj}
\textbf{Strategy.} On a line there is one direction to enforce; on a plane
with an orthonormal basis the projection splits into separate components
(Proposition~\ref{th:proj-onb}). Compute each piece, then check that the
residual kills the subspace.

\textbf{Solution.} Take $\mathbf a=(2,0,1)$ and $\mathbf b=(1,2,3)$.
Projection onto the line through $\mathbf a$:
\[
\mathbf p=\frac{\mathbf a^{T}\mathbf b}{\mathbf a^{T}\mathbf a}\,\mathbf a
=\frac{2+0+3}{4+0+1}\,\mathbf a=\frac55\,(2,0,1)=(2,0,1).
\]
Check: $\mathbf b-\mathbf p=(-1,2,2)$ and
$\mathbf a\cdot(\mathbf b-\mathbf p)=-2+0+2=0$. \checkmark

Now project the same $\mathbf b$ onto
$W=\operatorname{span}\{(1,1,0),(0,0,1)\}$. The spanning vectors are already
orthogonal; normalize to $\mathbf q_1=(1,1,0)/\sqrt2$, $\mathbf q_2=(0,0,1)$:
\[
\mathbf p_W=(\mathbf q_1^{T}\mathbf b)\mathbf q_1+(\mathbf q_2^{T}\mathbf b)\mathbf q_2
=\frac{3}{\sqrt2}\,\frac{(1,1,0)}{\sqrt2}+3(0,0,1)
=\bigl(\tfrac32,\tfrac32,3\bigr).
\]
Check: $\mathbf b-\mathbf p_W=(-\tfrac12,\tfrac12,0)$, which is orthogonal to
$\mathbf q_1$ ($-\tfrac12+\tfrac12=0$) and to $\mathbf q_2$. \checkmark
Notice $P_{\text{line}}$ was rank 1; the two pieces of $\mathbf p_W$ were
computed independently --- impossible with a skew basis.
\end{example}

\begin{figure}[htbp]
\centering
\begin{tikzpicture}[scale=1.0, >=stealth]
  % plane as parallelogram
  \fill[accentsoft!25] (-0.5,-0.3)--(4.3,-0.3)--(5.3,1.3)--(0.5,1.3)--cycle;
  \draw[gray] (-0.5,-0.3)--(4.3,-0.3)--(5.3,1.3)--(0.5,1.3)--cycle;
  \node[gray] at (4.9,0.6) {$W$};
  \coordinate (b) at (2.4,3);
  \coordinate (p) at (2.3,0.62);
  \draw[->, ink, very thick] (1.9,0.35)--(b) node[above]{$\mathbf b$};
  \draw[->, good, very thick] (1.9,0.35)--(p);
  \draw[->, amber, thick] (p)--(b) node[midway,right]{$\mathbf e$};
  \draw[accent, thick] (1.2,0.15)--(3.4,1.0) node[right]{$\mathbf u_1$};
  \draw[accent, thick] (1.5,0.85)--(3.8,0.1) node[right]{$\mathbf u_2$};
  \node[align=left] at (0.4,2.2) {\small $\mathbf e$ perpendicular to\\ \small every direction in $W$};
\end{tikzpicture}

\caption{Projecting $\mathbf b$ onto a plane $W$. Notice the residual
$\mathbf e$ is perpendicular to \emph{every} direction in $W$, not just the
basis vectors. Regenerate: \ttfamily figures/ch05\_projection\_plane.tex.}
\label{fig:proj-plane}
\end{figure}

The code checks the general formula end-to-end on a 4D example:

\begin{lstlisting}[caption={\ttfamily code/ch05\_projection\_idempotent.py ---
builds $P=A(A^{T}A)^{-1}A^{T}$ on a $4\times2$ basis, projects, and verifies
every defining property.}]
import numpy as np
np.set_printoptions(precision=4, suppress=True)

A = np.array([[1., 0.],
              [0., 1.],
              [1., 1.],
              [1., -1.]])
b = np.array([2., 1., 3., -1.])

P = A @ np.linalg.inv(A.T @ A) @ A.T
p = P @ b
print("P =")
print(P)
print("p = P @ b =", p)
print("in col(A)?", np.allclose(p, A @ np.linalg.solve(A.T @ A, A.T @ b)))
print("residual b - p =", b - p)
print("residual perpendicular to col(A)?", np.allclose(A.T @ (b - p), 0))
print("P^2 = P:", np.allclose(P @ P, P))
print("P^T = P:", np.allclose(P.T, P))
print("||b||^2 = ||p||^2 + ||b-p||^2:",
      np.isclose(b @ b, p @ p + (b - p) @ (b - p)))
\end{lstlisting}

\begin{lstlisting}[caption={Output of the projection check. Notice all three
signatures of an orthogonal projection at once: lands in $\operatorname{col}(A)$,
perpendicular residual, and $P^2=P$, $P^{T}=P$.}]
P =
[[ 0.3333  0.      0.3333  0.3333]
 [ 0.      0.3333  0.3333 -0.3333]
 [ 0.3333  0.3333  0.6667  0.    ]
 [ 0.3333 -0.3333  0.      0.6667]]
p = P @ b = [ 1.3333  1.6667  3.     -0.3333]
in col(A)? True
residual b - p = [ 0.6667 -0.6667  0.     -0.6667]
residual perpendicular to col(A)? True
P^2 = P: True
P^T = P: True
||b||^2 = ||p||^2 + ||b-p||^2: True
\end{lstlisting}

\section{Gram--Schmidt and the QR factorization}\label{sec:orth-gs}

Any basis can be upgraded to an orthonormal one by successive subtraction:
take each vector, remove the parts already claimed by earlier
orthonormal vectors, and normalize what is left.

The rule in one line: each new vector gives up the parts that earlier
orthonormal vectors already own. Formally, $\mathbf u_1=\mathbf a_1$, then
$\mathbf u_j=\mathbf a_j-\sum_{i<j}\operatorname{proj}_{\mathbf u_i}\mathbf a_j$,
and $\mathbf q_j=\mathbf u_j/\|\mathbf u_j\|$. The projection formula guarantees
$\mathbf u_j\perp\mathbf u_i$ for all $i<j$.

\begin{example}[Gram--Schmidt by hand]\label{ex:gs}
\textbf{Strategy.} Work left to right; after each step verify
$\mathbf u_j$ is orthogonal to every $\mathbf u_i$ before it. Fractions are
easier to check than radicals --- keep vectors unnormalized until the last
line.

\textbf{Solution.} Take
$\mathbf a_1=(1,1,0)$, $\mathbf a_2=(1,0,1)$, $\mathbf a_3=(0,1,1)$.
\[
\mathbf u_1=(1,1,0),\qquad
\mathbf u_2=\mathbf a_2-\frac{\mathbf a_2\cdot\mathbf u_1}{\mathbf u_1\cdot\mathbf u_1}\,\mathbf u_1
=(1,0,1)-\tfrac12(1,1,0)=\bigl(\tfrac12,-\tfrac12,1\bigr).
\]
Check $\mathbf u_1\cdot\mathbf u_2=\tfrac12-\tfrac12=0$. \checkmark\ For the
third vector, subtract \emph{both} projections:
\[
\mathbf u_3=\mathbf a_3-\frac{\mathbf a_3\cdot\mathbf u_1}{2}\,\mathbf u_1
-\frac{\mathbf a_3\cdot\mathbf u_2}{\mathbf u_2\cdot\mathbf u_2}\,\mathbf u_2
=(0,1,1)-\tfrac12(1,1,0)-\tfrac{1/2}{3/2}\,\bigl(\tfrac12,-\tfrac12,1\bigr)
=\bigl(-\tfrac23,\tfrac23,\tfrac23\bigr).
\]
Check: $\mathbf u_3\cdot\mathbf u_1=-\tfrac23+\tfrac23=0$ and
$\mathbf u_3\cdot\mathbf u_2=-\tfrac13-\tfrac13+\tfrac23=0$. \checkmark\
Normalizing ($\|\mathbf u_1\|=\sqrt2$, $\|\mathbf u_2\|=\sqrt{3/2}$,
$\|\mathbf u_3\|=2/\sqrt3$):
\[
\mathbf q_1=\tfrac{1}{\sqrt2}(1,1,0),\quad
\mathbf q_2=\tfrac{1}{\sqrt6}(1,-1,2),\quad
\mathbf q_3=\tfrac{1}{\sqrt3}(-1,1,1).
\]
\end{example}

\begin{figure}[htbp]
\centering
\begin{tikzpicture}[scale=0.9, >=stealth]
  % step 0
  \begin{scope}
    \draw[->,gray] (-0.3,0)--(2.3,0); \draw[->,gray] (0,-0.3)--(0,2.3);
    \draw[->,accent,very thick] (0,0)--(1.8,0.5) node[right]{$\mathbf a_1$};
    \draw[->,accent,very thick] (0,0)--(1.1,1.9) node[above]{$\mathbf a_2$};
    \node at (1,-0.6) {\small input: $\mathbf a_1,\mathbf a_2$};
  \end{scope}
  % step 1
  \begin{scope}[xshift=3.6cm]
    \draw[->,gray] (-0.3,0)--(2.3,0); \draw[->,gray] (0,-0.3)--(0,2.3);
    \draw[->,accent,very thick] (0,0)--(1.8,0.5);
    \draw[->,good,very thick] (0,0)--(0.95,0.26) node[below right=-2pt]{\small $\mathbf u_1$};
    \draw[->,amber,very thick] (0.95,0.26)--(1.1,1.9);
    \draw[dashed,gray] (0,0)--(1.1,1.9);
    \node at (1,-0.6) {\small subtract $\mathbf a_2$'s $\mathbf u_1$-part};
  \end{scope}
  % step 2
  \begin{scope}[xshift=7.2cm]
    \draw[->,gray] (-0.3,0)--(2.3,0); \draw[->,gray] (0,-0.3)--(0,2.3);
    \draw[->,good,very thick] (0,0)--(0.95,0.26) node[below right=-2pt]{\small $\mathbf u_1$};
    \draw[->,good,very thick] (0,0)--(0.15,1.63) node[above left=-2pt]{\small $\mathbf u_2\perp\mathbf u_1$};
    \draw[dashed,gray] (0.95,0.26)--(0.15,1.63);
    \node at (1,-0.6) {\small residual is $\mathbf u_2$};
  \end{scope}
  % step 3
  \begin{scope}[xshift=10.8cm]
    \draw[->,gray] (-0.3,0)--(2.3,0); \draw[->,gray] (0,-0.3)--(0,2.3);
    \draw[->,accent,very thick] (0,0)--(0.95,0.26) node[below right=-2pt]{\small $\mathbf q_1$};
    \draw[->,accent,very thick] (0,0)--(0.15,1.63) node[above left=-2pt]{\small $\mathbf q_2$};
    \node at (1,-0.6) {\small normalize: $\mathbf q_i=\mathbf u_i/\|\mathbf u_i\|$};
  \end{scope}
\end{tikzpicture}

\caption{Gram--Schmidt on two vectors. Notice the algorithm never creates a
new direction --- it strips each input of the directions already spoken
for, and the residual that survives is automatically perpendicular.
Regenerate: \ttfamily figures/ch05\_gs\_steps.tex.}
\label{fig:gs-steps}
\end{figure}

Collect the unit vectors as columns of $Q$ and reverse the bookkeeping:
since $\mathbf a_j=\sum_{i\le j}r_{ij}\mathbf q_i$ with
$r_{ij}=\mathbf q_i^{T}\mathbf a_j$ and $r_{jj}=\|\mathbf u_j\|$,

\begin{theorem}[QR factorization]\index{QR factorization}
Every $A\in\mathbb R^{m\times n}$ with linearly independent columns factors
as
\[
A=QR,
\]
where $Q\in\mathbb R^{m\times n}$ has orthonormal columns and
$R\in\mathbb R^{n\times n}$ is upper triangular with positive diagonal.
\end{theorem}

The factorization is a solver's dream: $Ax=b$ becomes $Rx=Q^{T}b$, and an
upper-triangular system costs one back-substitution.

\begin{example}[QR and solving $Ax=b$]\label{ex:qr}
\textbf{Strategy.} Read $Q$ off Example~\ref{ex:gs}, then get $R=Q^{T}A$
(the entries are exactly the coefficients discarded during Gram--Schmidt).
Solve by projecting $\mathbf b$ then back-substituting.

\textbf{Solution.} With the three vectors of Example~\ref{ex:gs},
\[
Q=\begin{bmatrix}1/\sqrt2&1/\sqrt6&-1/\sqrt3\\
1/\sqrt2&-1/\sqrt6&1/\sqrt3\\0&2/\sqrt6&1/\sqrt3\end{bmatrix},
\qquad
R=\begin{bmatrix}\sqrt2&1/\sqrt2&1/\sqrt2\\
0&\sqrt{3/2}&1/\sqrt6\\0&0&2/\sqrt3\end{bmatrix},
\]
and you can multiply $QR$ to recover $A$. For $\mathbf b=(3,2,1)$:
$y=Q^{T}\mathbf b=(5/\sqrt2,\ \sqrt{3/2},\ 0)$, and back-substitution on
$Rx=y$ gives $x_3=0$, $x_2=1$, $x_1=2$ --- so $x=(2,1,0)$, and indeed
$2\mathbf a_1+\mathbf a_2=(3,2,1)=\mathbf b$.
\end{example}

\begin{lstlisting}[caption={\ttfamily code/ch05\_qr\_solve.py --- factors the
same $3\times3$ with \ttfamily np.linalg.qr and solves $Ax=b$ via
$y=Q^{T}b$ and back-substitution. Sign convention: \ttfamily numpy returns
negative diagonal $R$ (both conventions are valid QR factorizations --- it
is $Q$'s sign that flips to compensate).}]
import numpy as np
np.set_printoptions(precision=4, suppress=True)

A = np.array([[1., 1., 0.],
              [1., 0., 1.],
              [0., 1., 1.]])
b = np.array([3., 2., 1.])

Q, R = np.linalg.qr(A)
print("Q =")
print(Q)
print("R =")
print(R)
print("Q^T Q = I:", np.allclose(Q.T @ Q, np.eye(3)))
print("QR = A:", np.allclose(Q @ R, A))

y = Q.T @ b
x = np.linalg.solve(R, y)          # triangular -> back substitution
print("x via QR =", x)
print("np.linalg.solve =", np.linalg.solve(A, b))
print("residual:", np.linalg.norm(A @ x - b))
\end{lstlisting}

\begin{lstlisting}[caption={Output of the QR solve. Notice the two
conventions agree on everything that matters: $Q^{T}Q=I$, $QR=A$, and the
same $x=(2,1,0)$ as the hand computation.}]
Q =
[[-0.7071  0.4082 -0.5774]
 [-0.7071 -0.4082  0.5774]
 [-0.      0.8165  0.5774]]
R =
[[-1.4142 -0.7071 -0.7071]
 [ 0.      1.2247  0.4082]
 [ 0.      0.      1.1547]]
Q^T Q = I: True
QR = A: True
x via QR = [ 2.  1. -0.]
np.linalg.solve = [2. 1. 0.]
residual: 6.753223014464259e-16
\end{lstlisting}

There is a trap hidden in the arithmetic of Gram--Schmidt itself. In the
\emph{classical} version, the coefficient for direction $i$ is computed
from the original $\mathbf a_j$; in the \emph{modified} version it is
computed from the partially-deflated vector --- same value in exact
arithmetic, very different in floating point.

\begin{lstlisting}[caption={\ttfamily code/ch05\_gs\_stability.py ---
classical vs.\ modified Gram--Schmidt on nearly-parallel columns
($\varepsilon=10^{-8}$). Measures orthogonality drift as
$\|Q^{T}Q-I\|_\infty$.}]
import numpy as np
np.set_printoptions(precision=2, suppress=True)
np.seterr(all="ignore")

def classical_gs(A):
    Q = np.zeros_like(A)
    for j in range(A.shape[1]):
        v = A[:, j].copy()
        for i in range(j):
            v -= (Q[:, i] @ A[:, j]) * Q[:, i]
        Q[:, j] = v / np.linalg.norm(v)
    return Q

def modified_gs(A):
    Q = np.zeros_like(A)
    V = A.copy()
    for j in range(A.shape[1]):
        Q[:, j] = V[:, j] / np.linalg.norm(V[:, j])
        for k in range(j + 1, A.shape[1]):
            V[:, k] -= (Q[:, j] @ V[:, k]) * Q[:, j]
    return Q

eps = 1e-8
A = np.array([[1., 1., 1.],
              [eps, 0., 0.],
              [0., eps, 0.],
              [0., 0., eps]])

Qc = classical_gs(A)
Qm = modified_gs(A)
print("classical  |Q^T Q - I| =", np.linalg.norm(Qc.T @ Qc - np.eye(3), ord=np.inf))
print(Qc.T @ Qc)
print("modified   |Q^T Q - I| =", np.linalg.norm(Qm.T @ Qm - np.eye(3), ord=np.inf))
print(Qm.T @ Qm)
Q_lib, _ = np.linalg.qr(A)
print("np.linalg.qr |Q^T Q - I| =",
      np.linalg.norm(Q_lib.T @ Q_lib - np.eye(3), ord=np.inf))
\end{lstlisting}

\begin{lstlisting}[caption={Output of the stability test. Notice classical
Gram--Schmidt loses \emph{half a unit} of orthogonality while modified loses
$10^{-8}$, and the library QR (Householder reflectors) is exact to machine
precision. This is not a toy gap --- it is why nobody ships classical
Gram--Schmidt.}]
classical  |Q^T Q - I| = 0.500000007071068
[[ 1.  -0.  -0. ]
 [-0.   1.   0.5]
 [-0.   0.5  1. ]]
modified   |Q^T Q - I| = 1.1153550716504105e-08
[[ 1. -0. -0.]
 [-0.  1.  0.]
 [-0.  0.  1.]]
np.linalg.qr |Q^T Q - I| = 7.726045597493973e-16
\end{lstlisting}

\begin{tripwire}[{Assuming classical and modified Gram--Schmidt are the same
algorithm}]
They are algebraically identical and numerically different. Classical GS
computes each projection coefficient against the original column, so
rounding errors in early $\mathbf q_i$ pollute everything downstream;
modified GS recomputes against the partially-deflated vector. On nearly
dependent input the classical version can return columns at $45^\circ$
angles, as the listing shows.
\end{tripwire}

\section{Orthogonal complements: the four subspaces pair up}\label{sec:orth-complements}

\begin{definition}[Orthogonal complement]\index{orthogonal complement}
The \emph{orthogonal complement} of a subspace $W\subseteq\mathbb R^n$ is
\[
W^\perp=\{\mathbf v\in\mathbb R^n:\langle\mathbf v,\mathbf w\rangle=0
\text{ for all }\mathbf w\in W\}.
\]
It is a subspace, $W\cap W^\perp=\{\mathbf0\}$, and
$\dim W+\dim W^\perp=n$.
\end{definition}

\begin{theorem}[Fundamental theorem, geometric form]\label{th:fund-orth}
For $A\in\mathbb R^{m\times n}$:
\[
\operatorname{row}(A)^{\perp}=\ker(A),\qquad
\operatorname{col}(A)^{\perp}=\ker(A^{T}),
\]
i.e.\ $\mathbb R^n=\operatorname{row}(A)\oplus\ker(A)$ and
$\mathbb R^m=\operatorname{col}(A)\oplus\ker(A^{T})$.
\end{theorem}

\begin{proof}
$\mathbf x\in\ker(A)$ iff $A\mathbf x=\mathbf0$ iff every row of $A$ is
orthogonal to $\mathbf x$ iff $\mathbf x\perp\operatorname{row}(A)$. Apply the
same argument to $A^{T}$.
\end{proof}

\begin{figure}[htbp]
\centering
\begin{tikzpicture}[scale=1.0, >=stealth]
  \draw[graybg!40, fill=graybg!25] (0,0) rectangle (4.6,3);
  \draw[dashed, gray] (2.3,0) -- (2.3,3);
  \node[align=center] at (1.15,2.6) {$\operatorname{row}(A)$\\ \small dim $r$};
  \node[align=center] at (3.45,2.6) {$\ker(A)$\\ \small dim $n-r$};
  \node[below] at (2.3,-0.15) {domain $\mathbb{R}^n$};
  \draw[<->,accent,thick] (1.15,0.9)--(3.45,0.9) node[midway,below]{\tiny $\perp$ complements};
  \draw[graybg!40, fill=graybg!25] (7.4,0) rectangle (12,3);
  \draw[dashed, gray] (7.4,1.5) -- (12,1.5);
  \node[align=center] at (9.7,2.25) {$\operatorname{col}(A)$,\quad dim $r$};
  \node[align=center] at (9.7,0.75) {$\ker(A^{T})$,\quad dim $m-r$};
  \node[below] at (9.7,-0.15) {codomain $\mathbb{R}^m$};
  \draw[<->,accent,thick] (8.6,1.9)--(10.8,1.0) node[midway,sloped,below=-1pt]{\tiny $\perp$};
  \draw[->, accent, very thick] (4.8,2.25) -- (7.2,2.25) node[midway, above]{$A$};
\end{tikzpicture}

\caption{The four subspaces as orthogonal pairs. Notice every vector
splits: its $\operatorname{row}(A)$ part is what $A$ ``sees,'' its
$\ker(A)$ part maps to $\mathbf0$. Regenerate:
\ttfamily figures/ch05\_complements.tex.}
\label{fig:complements}
\end{figure}

\begin{figure}[htbp]
\centering
\begin{tikzpicture}[scale=0.9, >=stealth]
  \begin{scope}
    \draw[->,gray] (-0.3,0)--(2.5,0); \draw[->,gray] (0,-0.3)--(0,2.5);
    \draw[->,accent,very thick] (0,0)--(2,0) node[right]{$\mathbf b_1$};
    \draw[->,accent,very thick] (0,0)--(0,2) node[above]{$\mathbf b_2$};
    \draw[dashed,gray] (2,0)--(2,2)--(0,2);
    \node at (1.1,-0.6) {\small orthogonal basis};
  \end{scope}
  \begin{scope}[xshift=5.5cm]
    \draw[->,gray] (-0.3,0)--(2.8,0); \draw[->,gray] (0,-0.3)--(0,2.5);
    \draw[->,amber,very thick] (0,0)--(2.3,0) node[right]{$\mathbf c_1$};
    \draw[->,amber,very thick] (0,0)--(1.5,2) node[above]{$\mathbf c_2$};
    \draw[dashed,gray] (2.3,0)--(3.8,2)--(1.5,2);
    \node at (1.4,-0.6) {\small skew basis, same plane};
  \end{scope}
\end{tikzpicture}

\caption{Orthogonal vs.\ skew bases of the same plane. Notice both span
the same space --- only the orthogonal one makes projection and
coordinates decouple. Regenerate: \ttfamily figures/ch05\_ortho\_vs\_skew.tex.}
\label{fig:ortho-vs-skew}
\end{figure}

\begin{exercise}[medium]\label{ex:orth-2}
Find the orthogonal complement of $W=\operatorname{span}\{(1,1,1)\}$ in
$\mathbb R^3$.
\begin{solution}
$W^\perp$ is the plane $x+y+z=0$, i.e.
$\operatorname{span}\{(1,-1,0),\,(1,0,-1)\}$. Dimension check:
$\dim W+\dim W^\perp=1+2=3$.
\end{solution}
\end{exercise}

\begin{exercise}[medium]\label{ex:orth-3}
For $A=\begin{bmatrix}1&0&1\\0&1&1\end{bmatrix}$, split
$\mathbf b=(1,1,0)$ into $\operatorname{row}(A)$ and $\ker(A)$ parts.
\begin{solution}
$\ker(A)$ is the line $t(1,1,-1)$. Projecting $\mathbf b$ onto it gives
$\frac{2}{3}(1,1,-1)$, so
$\mathbf b=\underbrace{\tfrac13(1,1,2)}_{\text{row part}}
+\underbrace{\tfrac13(2,2,-2)}_{\text{null part}}$; check
$A\cdot\frac13(2,2,-2)=\frac13(0,0)=\mathbf0$. \checkmark
\end{solution}
\end{exercise}

\begin{exercise}[hard]\label{ex:orth-4}
Let $P$ be the projection matrix onto a subspace $W$. Prove $I-P$ is the
projection matrix onto $W^\perp$.
\begin{solution}
$(I-P)^2=I-2P+P^2=I-P$ and $(I-P)^{T}=I-P$, so $I-P$ is a projection.
Its column space is $\{\mathbf v-P\mathbf v\}=\{\text{residuals}\}\subseteq W^\perp$;
for any $\mathbf w\in W^\perp$, $P\mathbf w=\mathbf0$ so
$(I-P)\mathbf w=\mathbf w$ and the image is all of $W^\perp$.
\end{solution}
\end{exercise}

\begin{example}[The pairing on one matrix]\label{ex:pairing}
For $A=\begin{bmatrix}1&2\\2&4\\3&6\end{bmatrix}$, exhibit
$\operatorname{col}(A)^{\perp}=\ker(A^{T})$ as \emph{the same line}
computed two ways.
\par\smallskip
\textbf{Strategy.} $\operatorname{col}(A)$ is a line through
$(1,2,3)$; its complement is a plane. $\ker(A^{T})$ is computed by
solving $A^{T}\mathbf y=\mathbf0$ --- one elimination.
\begin{solution}
$A^{T}=\begin{bmatrix}1&2&3\\2&4&6\end{bmatrix}$, RREF
$\begin{bmatrix}1&2&3\\0&0&0\end{bmatrix}$, so $\ker(A^{T})=
\operatorname{span}\{(-2,1,0),(-3,0,1)\}$ --- a plane. And
$(1,2,3)^{\perp}$ is exactly $x+2y+3z=0$: the same plane.
Basis-check: $(-2,1,0)\cdot(1,2,3)=0$ and
$(-3,0,1)\cdot(1,2,3)=0$. \checkmark\ Dimensions close the loop
too: the column line plus the kernel plane tile all of
$\mathbb{R}^{3}$ ($1+2=3$).
\end{solution}
\end{example}

\section{Hand derivation}\label{sec:handderiv5}

Derive the projection formula from scratch. Drop a perpendicular from
$\mathbf b$ to the line through $\mathbf a$. The foot $\mathbf p$ must satisfy
two conditions: $\mathbf p$ lies on the line, so $\mathbf p=\hat b\,\mathbf a$
for some scalar $\hat b$; and the residual $\mathbf e=\mathbf b-\mathbf p$ is
orthogonal to the line, so $\langle\mathbf a,\mathbf b-\hat b\mathbf a\rangle=0$.
Expanding,
$\langle\mathbf a,\mathbf b\rangle=\hat b\,\langle\mathbf a,\mathbf a\rangle$,
hence
\[
\hat b=\frac{\mathbf a^{T}\mathbf b}{\mathbf a^{T}\mathbf a},\qquad
\mathbf p=\frac{\mathbf a^{T}\mathbf b}{\mathbf a^{T}\mathbf a}\,\mathbf a .
\]
Now promote it to a subspace. With columns of $A$ spanning $W$, write
$\mathbf p=A\mathbf x$ and require the residual orthogonal to every column:
$A^{T}(\mathbf b-A\mathbf x)=\mathbf0$, i.e.\ $A^{T}A\mathbf x=A^{T}\mathbf b$.
Independent columns make $A^{T}A$ invertible, so
$\mathbf x=(A^{T}A)^{-1}A^{T}\mathbf b$ and
$P=A(A^{T}A)^{-1}A^{T}$. Finally check the algebra the geometry
promised: $P^2=A(A^{T}A)^{-1}A^{T}A(A^{T}A)^{-1}A^{T}=P$ since
$A^{T}A$ cancels, and $P^{T}=P$ by expanding the transpose.

\section{Common mistakes}

\begin{tripwire}[Projecting onto the wrong subspace]
$\operatorname{proj}_W\mathbf b$ answers ``closest point in $W$,'' not
``closest point \emph{to} $W$.'' In least squares you project $\mathbf b$
onto $\operatorname{col}(A)$, not onto the rows --- swapping them is the
most common conceptual bug in the next chapter.
\end{tripwire}

\begin{tripwire}[{Forgetting to normalize before projecting onto a direction}]
$\operatorname{proj}_{\mathbf a}\mathbf b
=\frac{\mathbf a^{T}\mathbf b}{\mathbf a^{T}\mathbf a}\mathbf a$ --- the
denominator is $\|\mathbf a\|^2$, not $\|\mathbf a\|$. If $\mathbf a$ is
already a unit vector they coincide, which is why orthonormal bases make
the formula look like ``just dot products.''
\end{tripwire}

\begin{tripwire}[Treating similar-looking projections as interchangeable]
Projecting with an orthonormal basis ($\mathbf p=\sum_{i=1}^{k}(\mathbf q_i^{T}\mathbf b)\mathbf q_i$)
and with a general basis ($P=A(A^{T}A)^{-1}A^{T}\mathbf b$) give the same
point only when the orthonormal vectors span the same $W$. The second
formula works without orthogonality but costs an inverse.
\end{tripwire}

\begin{tripwire}[{Assuming $P^2=P$ and $P^{T}=P$ are optional properties}]
They are not properties you must verify case-by-case; they \emph{define}
orthogonal projection. If a ``projection'' you computed fails either one,
it is wrong --- or it was an oblique (non-orthogonal) projection to begin
with.
\end{tripwire}

\section{Summary}

\begin{center}
\begin{tabularx}{\textwidth}{l>{\raggedright\arraybackslash}X}
\toprule
\textbf{Object} & \textbf{What you need to know}\\
\midrule
Inner product & $\langle\mathbf u,\mathbf v\rangle=\mathbf u^{T}\mathbf v$;
gives norm $\|\mathbf v\|=\sqrt{\langle\mathbf v,\mathbf v\rangle}$ and
$\cos\theta=\langle\mathbf u,\mathbf v\rangle/\|\mathbf u\|\|\mathbf v\|$ \\
Orthonormal basis & $\mathbf q_i^{T}\mathbf q_j=\delta_{ij}$; coordinates are
dot products $c_j=\mathbf q_j^{T}\mathbf v$ \\
Orthogonal matrix & $Q^{T}Q=I$, i.e.\ $Q^{-1}=Q^{T}$; preserves lengths
and angles; $\det=\pm1$ \\
Projection onto line & $\mathbf p=\frac{\mathbf a^{T}\mathbf b}{\mathbf a^{T}\mathbf a}\mathbf a$;
$P=\mathbf a\mathbf a^{T}/\mathbf a^{T}\mathbf a$ \\
Projection onto subspace & orthonormal basis:
$\mathbf p=\sum_{i=1}^{k}(\mathbf q_i^{T}\mathbf b)\mathbf q_i$; general basis:
$P=A(A^{T}A)^{-1}A^{T}$; always $P^2=P$, $P^{T}=P$ \\
Gram--Schmidt & $\mathbf u_j=\mathbf a_j-\sum_{i<j}\operatorname{proj}_{\mathbf u_i}\mathbf a_j$,
$\mathbf q_j=\mathbf u_j/\|\mathbf u_j\|$; use the modified version \\
QR & $A=QR$: orthonormal $Q$, upper-triangular $R$; solves
$Ax=b$ as $Rx=Q^{T}b$ \\
Orthogonal complement & $\operatorname{row}(A)^{\perp}=\ker(A)$,
$\operatorname{col}(A)^{\perp}=\ker(A^{T})$; every space splits
$W\oplus W^\perp=\mathbb R^n$ \\
\bottomrule
\end{tabularx}
\end{center}
\begin{rulebox}[{Watch instead}]
\begin{enumerate}
\item \emph{Dot products and duality} --- 3Blue1Brown, Essence of
Linear Algebra ch.~9, 14:11, \url{https://youtu.be/LyGKycYT2v0} ---
the inner-product machinery of \S\ref{sec:orth-inner}
(accessed 2026-10-11).
\item \emph{Projections onto Subspaces} --- MIT OpenCourseWare, 18.06
Lecture 15, 48:51, \url{https://youtu.be/Y_Ac6KiQ1t0} ---
\S\ref{sec:orth-proj} at full length, same notation
(accessed 2026-10-11).
\end{enumerate}
\end{rulebox}

\section*{Further reading}
\addcontentsline{toc}{section}{Further reading}

\begin{itemize}
\item G.~Strang, \emph{Introduction to Linear Algebra}, 6th ed., ch.~4
(textbook level) --- orthogonality, projections, and Gram--Schmidt as
this chapter teaches them.
\item MIT OpenCourseWare, \emph{18.06SC Linear Algebra}, Unit II ---
recitations on projections and QR, at OCW (course ``18.06SC'', Unit
II page), \url{https://ocw.mit.edu}. Verified 2026-10-11.
\item S.~Boyd and L.~Vandenberghe, \emph{VMLS}, chs.~3, 5 (free PDF) ---
distance, angle, and independence in the data-first register.\newline
\url{https://web.stanford.edu/~boyd/vmls/vmls.pdf}
Verified 2026-10-11.
\item Margalit and Rabinoff, \emph{Interactive Linear Algebra},
``Orthogonality'' --- dot products through Gram--Schmidt,
hyperlinked.\newline
\url{https://textbooks.math.gatech.edu/ila/}
Verified 2026-10-11.
\item M.~Deisenroth et al., \emph{Mathematics for Machine Learning},
ch.~3 (free PDF) --- inner products and orthogonal projections for ML.
\newline
\url{https://mml-book.github.io/book/mml-book.pdf}
Verified 2026-10-11.
\item J.~Hefferon, \emph{Linear Algebra}, ch.~VI topic ``Gram--Schmidt''
(free textbook) --- the rigorous version of the same
orthogonalization.\newline
\url{https://hefferon.net/linearalgebra/index.html}
Verified 2026-10-11.
\end{itemize}

\section*{Exercises}

\begin{exercise}[easy]
Compute $\operatorname{proj}_{\mathbf a}\mathbf b$ for
$\mathbf a=(1,2,2)$, $\mathbf b=(4,-1,1)$, and verify the residual is
orthogonal to $\mathbf a$.
\end{exercise}

\begin{exercise}[easy]
Is $Q=\frac{1}{5}\begin{bmatrix}3&4\\4&-3\end{bmatrix}$ orthogonal? What is
its determinant, and which kind of map is it?
\end{exercise}

\begin{exercise}[easy]
Run Gram--Schmidt by hand on $(1,0),(1,1)$, then write the QR of
$A=\begin{bmatrix}1&1\\0&1\end{bmatrix}$.
\end{exercise}

\begin{exercise}[easy]
Show that if $W$ is a subspace, so is $W^\perp$ (check the three
subspace properties from ch.~3).
\end{exercise}

\begin{exercise}[medium]
Prove the projection $\mathbf p=\operatorname{proj}_W\mathbf b$ is the unique
closest point of $W$ to $\mathbf b$: for any $\mathbf w\in W$,
$\|\mathbf b-\mathbf w\|^2=\|\mathbf b-\mathbf p\|^2+\|\mathbf p-\mathbf w\|^2$.
(Hint: Pythagoras --- where do the vectors live?)
\end{exercise}

\begin{exercise}[medium]
A projection matrix has $P^{T}=P$, $P^2=P$. Show its eigenvalues (ch.~7
preview) can only be 0 or 1: if $P\mathbf v=\lambda\mathbf v$ then
$\lambda^2=\lambda$.
\end{exercise}

\begin{exercise}[medium]
Compute the QR of $A=\begin{bmatrix}1&1\\1&0\\0&1\end{bmatrix}$ by hand and
use it to project $\mathbf b=(1,1,1)$ onto $\operatorname{col}(A)$. Verify
$A^{T}(\mathbf b-\mathbf p)=\mathbf0$.
\end{exercise}

\begin{exercise}[medium]
In code (\texttt{numpy}), generate a random $6\times3$ matrix $A$, compute
$P=A(A^{T}A)^{-1}A^{T}$, and verify $P^2=P$, $P^{T}=P$, and
$\operatorname{rank}(P)=3$. What do the eigenvalues of $P$ look like?
\end{exercise}

\begin{exercise}[medium]
For $A=\begin{bmatrix}1&2\\2&4\\3&6\end{bmatrix}$ (rank 1), write the orthogonal
decomposition of $\mathbf b=(1,1,1)$ into $\operatorname{col}(A)$ and
$\ker(A^{T})$ parts. Which subspace holds most of $\mathbf b$?
\end{exercise}

\begin{exercise}[hard]
Prove that a square matrix $P$ is an orthogonal projection
($P^2=P$, $P^{T}=P$) onto $\operatorname{col}(P)$ --- i.e.\ the two algebraic
conditions characterize projections completely.
\end{exercise}

\begin{exercise}[hard]
Show that $(W^\perp)^\perp=W$ for any subspace $W\subseteq\mathbb R^n$, and
use it to prove $\operatorname{row}(A)=\ker(A)^\perp$.
\end{exercise}

\begin{exercise}[hard]
Implement modified Gram--Schmidt yourself and compare against
\texttt{np.linalg.qr} on a random $8\times5$ matrix: report
$\|Q^{T}Q-I\|_\infty$ and $\|A-QR\|_\infty$ for both, plus the residual of
$Rx=Q^{T}b$ solving $Ax=b$. Explain any sign differences in $Q$ and $R$.
(Reference implementation: \texttt{code/ch05\_gs\_stability.py}.)
\end{exercise}
\begin{exercise}[easy]
Compute $\cos\theta$ between $\mathbf{u}=(1,0)$ and
$\mathbf{v}=(1,1)$, then between $\mathbf u$ and $-\mathbf v$. What
does the sign flip say geometrically?
\end{exercise}
\begin{exercise}[medium]
Let $\mathbf{q}_1=(1,1,0)/\sqrt2$ and $\mathbf{q}_2=(1,-1,0)/\sqrt2$.
Extend to an orthonormal basis of $\mathbb{R}^3$ --- find
$\mathbf{q}_3$ by inspection, then confirm by Gram--Schmidt on
$(1,0,0)$.
\end{exercise}
\begin{exercise}[hard]
Prove that orthogonal nonzero vectors are automatically independent,
then show the converse fails: exhibit independent vectors whose
Gram--Schmidt output differs from the input. What does the difference
measure?
\end{exercise}

\chapternotes{%
Projection onto a line and the Gram--Schmidt--QR pipeline follow
Strang's ordering (orthogonalization before least squares)
\cite{ila,ocw1806sc}. The classical-versus-modified stability gap and
Householder-style reflection QR are treated numerically in MML ch.~3
\cite{mml} and practically in VMLS ch.~11 \cite{vmls}. Orthogonal
complements of the four subspaces are the geometric half of the
fundamental theorem; the dimension half was the rank--nullity theorem
of ch.~3 \cite{axler,ila}.}
